\subsection{Software workflow and outputs}
\label{ch:code}

Once the inputs satisfy the data contract, the software must carry them
through the same sequence as the derivation: construct the transport block,
form the spatial Jacobian, solve the adjoint, and write results with their
diagnostics and provenance. The public interface exposes that sequence without
requiring an application to modify the solver.

\subsubsection{Public interface}

The package exposes the data-to-derivative sequence through a small typed API.
The same calls validate inputs, build the model selected in the configuration,
and write results with their provenance:
\begin{lstlisting}[language=Julia]
using TransportNetworkWelfare

project = load_project("examples/grid_multimodal/config.toml")
validate(project)
model = build_model(project)
welfare = welfare_effects(model)
components = decompose_welfare(model)
sensitivity = sensitivity_path(model, :eta, [0.9, 1.099, 1.25])
write_results(components, "examples/grid_multimodal/output/guide-api/"; project)
\end{lstlisting}

The configuration determines the spatial, modal, congestion, and policy
types. Users adapting a dataset should change CSV and TOML inputs rather than
editing the solver.

The CLI calls the same public functions:
\begin{lstlisting}[language=bash]
julia --project=. bin/tnw.jl validate \
  examples/grid_multimodal/config.toml
julia --project=. bin/tnw.jl analyze \
  examples/grid_multimodal/config.toml
julia --project=. bin/tnw.jl decompose \
  examples/grid_multimodal/config.toml
julia --project=. bin/tnw.jl sensitivity \
  examples/grid_multimodal/config.toml
\end{lstlisting}

\subsubsection{Project isolation and package updates}

Section \ref{ch:quickstart} gives the initial setup and run commands. Keep the
resulting empirical project outside the package checkout so application inputs
and outputs remain separate from package updates. Before updating, preserve
the package commit and manifest used for archived results. Then fast-forward
the package and rerun its tests:
\begin{lstlisting}[language=bash]
git rev-parse HEAD
git switch main
git pull --ff-only
julia --project=. -e 'using Pkg; Pkg.instantiate(); Pkg.test()'
\end{lstlisting}
Revalidate and rerun the external project after any package, configuration, or
data change. A changed output belongs to a new result vintage; retain the
earlier commit, inputs, and manifest when earlier results must remain
reproducible.

\subsubsection{Implementation boundary}

Users normally interact with the API or CLI rather than individual source
modules. The implementation follows the derivation's dependency order:
validate data, reconstruct routes and modes, construct the selected spatial
system, eliminate congestion feedback, and project each policy onto welfare.
Appendices \ref{app:crosswalk} and \ref{app:algorithm} map these operations to
source files and functions. The tracked theory--code audit records the
line-level correspondence for contributors.

\subsubsection{Build stages}

Figure \ref{fig:guide-pipeline} in the quick start summarizes the calculation.
Each arrow passes a checked result to the next stage, beginning with
model-ready inputs and ending with archived results.

For either spatial model, \code{build\_model} begins by loading and validating
the project and network, identifying the positive-flow edge--mode pairs, and
reconstructing the route kernel and bilateral flow matrix. It uses these
objects to form the flexible and fixed route derivatives. The
economic-geography or urban builder then supplies the no-congestion spatial
Jacobian, cost loading, and welfare row.

The remaining stages construct the modal and congestion maps for each
decomposition case and eliminate congestion quantities through the Schur
complement. After checking the route, transport, and equilibrium condition
numbers, the builder exposes the equilibrium system used for the welfare and
decomposition calculations.

\code{analyze} uses a welfare-only builder. It constructs the full flexible
case but omits the dense fixed-route incidence matrix. \code{decompose}
constructs all decomposition cases. Omitting fixed-route incidence from ordinary analysis
reduces memory use because that matrix scales with origins, destinations, and
edges.

\subsubsection{Policy derivatives}

For all directed policy arcs, the code constructs a sparse selector matrix.
The realized forcing changes the aggregate cost on the selected edge. The
primitive forcing first passes the edge--mode shock through the modal and route
responses. It then incorporates edge or terminal congestion, aggregates the
result to the directed edge, and loads the resulting cost change into the
spatial residual.

The adjoint gain function solves the transposed Jacobian and multiplies the
forcing matrix. It returns one elasticity per policy arc. The decomposition
then applies inverse-gap and structured-update identities to the same forcing.

\subsubsection{Outputs and manifests}

Analysis writes directed results to \code{welfare\_directed.csv} and, when
requested, reciprocal-link results to \code{welfare\_physical.csv}. The
\code{decompose} command adds the corresponding decomposition tables. Each run
also writes \code{diagnostics.json} and \code{run\_manifest.json}.

The manifest records the package version and Git state, Julia and BLAS
environment, configuration hash, input hashes, model variants, parameters,
transformations, output hashes, and verification status. JSON serialization is
strict: unsupported values and colliding stringified keys are errors rather
than silently converted strings.

\begin{longtable}{L{0.28\textwidth}L{0.39\textwidth}L{0.24\textwidth}}
\toprule
Field & Meaning & Unit or sign \\
\midrule
\endhead
\code{physical\_link\_id} &
Identifier shared by the reciprocal directed policies. &
Identifier. \\
\code{hulten} &
Traditional approach: baseline policy-mode value traffic. &
Positive welfare elasticity. \\
\code{realized\_F} &
Full-model response to a reduction in the cost users face. &
Benefit elasticity; may be zero. \\
\code{primitive\_F} &
Extended approach: full-model response to the primitive infrastructure cost. &
Headline benefit elasticity. \\
\code{chi\_effective} &
Ratio \(\code{primitive\_F}/\code{realized\_F}\). &
Missing when the denominator is numerically zero. \\
\code{primitive\_externality}\\
\code{primitive\_propagation} &
Common externality scale and link-specific equilibrium-propagation terms in
the signed Hulten gap. &
Signed welfare elasticities. \\
\code{primitive\_edge}\\
\code{primitive\_terminal}\\
\code{primitive\_pass\_through} &
Edge congestion, terminal congestion, and primitive-to-realized cost
contributions. &
Signed welfare elasticities. \\
\code{identity\_residual\_*} and \code{channel\_residual\_*} &
Reconstruction errors for the closure ladder and analytical-component identities. &
Absolute numerical residuals; target below \(10^{-10}\). \\
\bottomrule
\end{longtable}

The \code{decomposition\_physical.csv} file sums opposite-direction
elasticities and signed components before calculating link-level ratios and
normalized multipliers. Ranks are derived from the selected welfare column;
they are not stored as primitive output fields. Join the physical-link file to
engineering costs using \code{physical\_link\_id}, and preserve the mapping
used for that merge. A link is ready for economic comparison only when the
run-level accounting, conditioning, and identity diagnostics pass.

\subsubsection{Computational cost}

The sparse welfare calculation is appropriate for many candidate edges on a
moderate network. The exact route decomposition is more demanding because it
constructs baseline origin-destination edge incidence. Before decomposition,
the validator reports an incidence-memory estimate. Large traffic-assignment
standard networks may be suitable for the welfare derivative but not the full
fixed-route decomposition on a laptop.

The package does not replace an ill-conditioned model with low-rank fitting,
ridge regularization, or dropped routes. Larger applications should exploit
sparsity or use a separately derived route comparison while retaining the
verification gates.

The public grid exposes the intermediate quantities and diagnostics used in the
calculation. Before applying the architecture to that baseline, the next
subsection states the numerical and reporting checks that a calculation must
pass.
